

\begin{document}
	
	% RESET all counters at the beginning
	\setcounter{section}{0}
	\setcounter{subsection}{0}
	\setcounter{subsubsection}{0}
	\setcounter{paragraph}{0}
	
	% Depth for numbering and TOC
	\setcounter{secnumdepth}{1}  % Number sections only
	% Part in TOC: footnotesize, bold, NO page break
	\makeatletter
	\renewcommand*\l@part[2]{%
		\ifnum \c@tocdepth >-2\relax
		\addpenalty{-\@highpenalty}%
		\addvspace{0.8em \@plus\p@}%
		{\leftskip 0em \relax
			\rightskip \@tocrmarg
			\parfillskip -\rightskip
			\parindent 0em \relax\@afterindenttrue
			\interlinepenalty\@M
			\leavevmode
			{\footnotesize\bfseries #1}\nobreak
			\leaders\hbox{$\m@th\mkern \@dotsep mu\hbox{}\mkern \@dotsep mu$}\hfill
			\nobreak\hb@xt@\@pnumwidth{\hss #2}\par}%
		\addvspace{0.2em \@plus\p@}%
		\nobreak
		\fi}
	\makeatother
	
	% Chapter in TOC: footnotesize, bold
	\renewcommand{\cftchapfont}{\footnotesize\bfseries}
	\renewcommand{\cftchappagefont}{\footnotesize\bfseries}
	\setlength{\cftbeforechapskip}{0.3em}
	
	% Only Chapters in TOC (no Sections/Subsections)
	\setcounter{tocdepth}{0}
	

	\begin{titlepage}
		\centering
		\vspace*{2cm}
	{\Huge\textbf{Fundamentale Fraktal-Geometrische Feldtheorie (FFGFT) oder T0-Theorie: Zeit-Masse-Dualität}}\\[1cm]
	{\Large Teil 8: Standardmodell aus GF(27), Leptonmassen und Schnittstellen}\\[2cm]
		{\large Johann Pascher}\\[1cm]
		{\large 2026}
		\vfill
	\end{titlepage}

	\frontmatter
	\pagestyle{fancy}
	\makeatletter
	\let\ps@plain\ps@fancy
	\let\ps@empty\ps@fancy
	\makeatother

	\mainmatter
	\pagestyle{fancy}
	\tableofcontents

	\chapter{Einleitung zu Band 8}
	% =============================================================================
% EINLEITUNG ZU BAND 8: STANDARDMODELL AUS GF(27), LEPTONMASSEN UND
% SCHNITTSTELLEN (Bd. 8 von 8)
% =============================================================================

\section*{Achter Band der Dokumentensammlung}

Band~8 umfasst die Dokumente 344 bis 389 aus der Zeit von September bis
Anfang Oktober 2026 und schließt die Sammlung auf dem aktuellen Stand des
Korpus ab. Er baut unmittelbar auf der Galois-Brücke aus Band~7 auf und
führt sie bis zu den Quantenzahlen des Standardmodells, zu den Leptonmassen
und zu den Schnittstellen mit anderen theoretischen Rahmen weiter.

\subsection*{Inhalt und Schwerpunkt}

Der erste Teil leitet aus der Struktur von $GF(27)^*$ die
Ladungsquantisierung, die Gell-Mann--Nishijima-Relation
$Q=I_3+Y/2$, die Generationenstruktur mit ihrer Kopplungsmatrix und die
$SU(2)_L$-Chiralität ab; ergänzt wird dies durch den Vergleich mit den
Furey-Fermionen, eine Einordnung in die Literatur und die Behandlung
Schwarzer Löcher unter der Zeit-Masse-Dualität. Der zweite Teil klärt
zunächst, wie die PDG-Vergleichswerte zustande kommen und welche
Modellannahmen in sie eingehen, und stellt dann Leptonleiter, Galois-Ansatz
und Koide-Relation gegenüber; es folgen die Koide-Amplitude $d/c=\sqrt2$, die
Higgs-Geometrie, die spontane Symmetriebrechung aus der
$T^4/\mathbb{Z}_3$-Geometrie und die Übersicht der Standardmodell-Teilchen
mit ihrem Belegstatus.

Der dritte Teil zeigt, warum Harmonik und algebraische Struktur dieselbe
Mathematik sind, und prüft diese Einsicht an Laborsystemen: Frequenzanalyse
und ihre Grenzen, HRV-Messung, Kagome-Gitter und die BAW-Ising-Maschine.
Hinzu kommen die Hodge-Theorie auf $T^4/\mathbb{Z}_3$, der Vergleich von
Ikosaeder-Träger und $D_4$-Gitter, die Darstellung des Fundaments in vier
Schichten und die Feldstruktur von Energie, Fluss und Geometrie. Der vierte
Teil ist den Leptonmassen gewidmet: $\theta=p_0=2/9$ als zwei Lesarten eines
Quotienten, der Weg von $p_0,p_1,p_2$ zu den Massen, vier Wege im Vergleich,
die Begründung des Ikosaeders, der Faktor $4/3$, das Matrixelement-Prinzip
und der Yukawa-Mechanismus als Folge von $\tilde T\cdot m=1$.

Der fünfte Teil sammelt die Schnittstellen: die Brückengleichungen zu
Observer Patch Holography, die Hierarchie $v/E_P$, eine Übersicht aller
Vergleichs- und Brückendokumente, Überlegungen zur Begutachtung, die
Spezifität des $D_4$-Trägers, die exakt aufsummierte Rekursion, den Stand
des Myon-$g-2$, die Massen ohne $v$, eine Kurzfassung der FFGFT nach der
Rechenprüfung, den Higgs-Vakuum-Weg zu $\xi$, die Lage von $\alpha$ in der
Geometrie, eine Übersicht der Abweichungen, die CMB-Temperatur in der
Grundform und die Darstellungen der Gravitation am Beispiel des
Parabelflugs.

\subsection*{Gliederung}

Die 46 Dokumente sind in fünf Teile gegliedert: Quantenzahlen aus der
Galois-Struktur (Dok.~344--350); Vergleichswerte, Leptonreste, Higgs und
Standardmodell-Teilchen (Dok.~351--357); Harmonik, Laborsysteme,
Hodge-Theorie und Felder (Dok.~358--366); Ikosaeder, Matrixelemente und
Leptonmassen (Dok.~367--373); Schnittstellen, Begutachtung und Grundform
(Dok.~374--389).

\subsection*{Zum Lesen der Dokumente}

Für einen raschen Überblick über den Gesamtstand eignet sich Dok.~384
(Kurzfassung), für die verbleibenden Abweichungen Dok.~387. Statusmarker
und Korrekturregister gelten wie in den vorangehenden Bänden; maßgeblich ist
jeweils die jüngste Fassung eines Ergebnisses im Korpus.

\subsection*{Stand}

Die Dokumente dieses Bandes tragen den Stand ihrer jeweiligen Fassung im
Korpus zum Oktober 2026.


	% --- part ---
	\part{Teil I --- Quantenzahlen aus der Galois-Struktur}
	\FFGFTdoc{Dok.~344 --- Furey-Fermionen in GALG und FFGFT: N-Operator, Tripotente und Neutrino-Vakuum-Identität}{../ch/344_Furey_GALG_FFGFT_De_ch}
	\FFGFTdoc{Dok.~345 --- FFGFT im Literaturvergleich: Was leistet der Galois-Ansatz, was nicht}{../ch/345_FFGFT_Literaturvergleich_De_ch}
	\FFGFTdoc{Dok.~346 --- Ladungsquantisierung aus GF$(27)^*$: Elektrische Ladung und Farbladung algebraisch erzwungen}{../ch/346_Ladungsquantisierung_Galois_De_ch}
	\FFGFTdoc{Dok.~347 --- Gell-Mann-Nishijima aus der Galois-Struktur: $Q = I_3 + Y/2$ algebraisch hergeleitet}{../ch/347_Gell_Mann_Nishijima_Galois_De_ch}
	\FFGFTdoc{Dok.~348 --- Generationenstruktur und Kopplungsmatrix aus GF$(27)$: Spur-Bilinearform, Orbit-Zuordnung und CKM-Komplexität}{../ch/348_Generationen_Orbits_De_ch}
	\FFGFTdoc{Dok.~349 --- SU$(2)_L$-Chiralität aus der Galois-Struktur: Su/Sd-Trennung und Cl$(6)$-Gradparität}{../ch/349_SU2L_Chiralitaet_Galois_De_ch}
	\FFGFTdoc{Dok.~350 --- Schwarze Löcher und Zeit-Masse-Dualität: Kovarianz von $G$, Verdampfungssequenz und das Neutrino als Endzustand}{../ch/350_Schwarze_Loecher_Dualitaet_De_ch}
	% --- part ---
	\part{Teil II --- Vergleichswerte, Leptonreste, Higgs und Standardmodell-Teilchen}
	\FFGFTdoc{Dok.~351 --- Modellabhängigkeit der PDG-Vergleichswerte}{../ch/351_PDG_Modellabhaengigkeit_De_ch}
	\FFGFTdoc{Dok.~352 --- Leptonreste: Leptonleiter, Galois und Koide im Vergleich}{../ch/352_Leptonreste_De_ch}
	\FFGFTdoc{Dok.~353 --- Koide-Amplitude $d/c = \sqrt{2}$ und die $T^4/Z_3$-Geometrie}{../ch/353_Koide_Amplitude_De_ch}
	\FFGFTdoc{Dok.~354 --- $H^0$ und $\tilde{T}\cdot m=1$: Neun Erkenntnisse zur Higgs-Geometrie}{../ch/354_Higgs_TM_Dualitaet_De_ch}
	\FFGFTdoc{Dok.~355 --- SSB $SU(2)_L\times U(1)_Y\to U(1)_\text{EM}$ aus $T^4/Z_3$-Geometrie}{../ch/355_SSB_Torus_De_ch}
	\FFGFTdoc{Dok.~356 --- Standardmodell-Teilchen aus GF$(27)^*$}{../ch/356_SM_Teilchen_GF27_De_ch}
	\FFGFTdoc{Dok.~357 --- Standardmodell-Quantenzahlen aus GF$(27)^*$ --- Belegstatus}{../ch/357_SM_Teilchen_GF27_Zusatz_De_ch}
	% --- part ---
	\part{Teil III --- Harmonik, Laborsysteme, Hodge-Theorie und Felder}
	\FFGFTdoc{Dok.~358 --- Harmonik und algebraische Struktur}{../ch/358_Harmonik_Algebra_De_ch}
	\FFGFTdoc{Dok.~359 --- KI-basierte Mustererkennung und algebraische Grenzen der Frequenzanalyse}{../ch/359_KI_Grenzen_Frequenzanalyse_De_ch}
	\FFGFTdoc{Dok.~360 --- Galois-informierte HRV-Analyse}{../ch/360_HRV_Galois_Experiment_De_ch}
	\FFGFTdoc{Dok.~361 --- Kagome-Gitter und FFGFT: Vier strukturelle Parallelen}{../ch/361_Kagome_FFGFT_Vergleich_De_ch}
	\FFGFTdoc{Dok.~362 --- Die BAW-Ising-Maschine als Resonanzrechner}{../ch/362_BAW_Ising_FFGFT_De_ch}
	\FFGFTdoc{Dok.~363 --- Hodge-Theorie auf T\textsuperscript{4}/Z\textsubscript{3} im Rahmen der FFGFT}{../ch/363_Hodge_T4Z3_De_ch}
	\FFGFTdoc{Dok.~364 --- Der Ikosaeder-Träger und das D\textsubscript{4}-Gitter}{../ch/364_Ikosaeder_D4_Vergleich_De_ch}
	\FFGFTdoc{Dok.~365 --- Das Fundament der FFGFT in vier Schichten}{../ch/365_Fundament_FFGFT_De_ch}
	\FFGFTdoc{Dok.~366 --- Felder in der FFGFT: Energie, Fluss und Geometrie}{../ch/366_Feldstruktur_FFGFT_De_ch}
	% --- part ---
	\part{Teil IV --- Ikosaeder, Matrixelemente und Leptonmassen}
	\FFGFTdoc{Dok.~367 --- Warum $\theta = p_0 = 2/9$ --- Wahrscheinlichkeit und Phase als zwei Lesarten eines Quotienten}{../ch/367_Phase_Wahrscheinlichkeit_De_ch}
	\FFGFTdoc{Dok.~368 --- Von $p_0,p_1,p_2$ zu Leptonmassen}{../ch/368_p_Werte_Lepton_Massen_De_ch}
	\FFGFTdoc{Dok.~369 --- Vier Wege zu den Leptonmassen}{../ch/369_Vier_Wege_Leptonmassen_De_ch}
	\FFGFTdoc{Dok.~370 --- Warum das Ikosaeder --- Drei, Fünf und das Nicht-Kommutieren}{../ch/370_Warum_Ikosaeder_De_ch}
	\FFGFTdoc{Dok.~371 --- Der Faktor 4/3 --- Kugelvolumen, Casimir-Operatoren und elektromagnetische Masse}{../ch/371_Faktor_Vier_Drittel_De_ch}
	\FFGFTdoc{Dok.~372 --- Das Matrixelement-Prinzip und seine Reichweite}{../ch/372_Matrixelement_Mechanismus_De_ch}
	\FFGFTdoc{Dok.~373 --- Der Yukawa-Mechanismus als Konsequenz von $\tilde T\cdot m=1$}{../ch/373_Yukawa_Wicklungszahlen_De_ch}
	% --- part ---
	\part{Teil V --- Schnittstellen, Begutachtung und Grundform}
	\FFGFTdoc{Dok.~374 --- Gemeinsamer Anker --- Brückengleichungen zu Observer Patch Holography}{../ch/374_Gemeinsamer_Anker_OPH_De_ch}
	\FFGFTdoc{Dok.~375 --- Die Hierarchie $v/E_P$ --- Zwei Herleitungen im Vergleich: FFGFT und Observer Patch Holography}{../ch/375_Hierarchie_vEP_Vergleich_De_ch}
	\FFGFTdoc{Dok.~376 --- Der Koeffizient 8$\pi$ und der Status von Issue 740}{../ch/376_Koeffizient_8pi_740_De_ch}
	\FFGFTdoc{Dok.~377 --- Schnittstellen der FFGFT zu anderen Rahmen}{../ch/377_Schnittstellen_FFGFT_De_ch}
	\FFGFTdoc{Dok.~378 --- Begutachtung unter heutigen Bedingungen}{../ch/378_Begutachtung_heute_De_ch}
	\FFGFTdoc{Dok.~379 --- Das akustische Plenum und das Photon}{../ch/379_Lien_Plenum_Photon_De_ch}
	\FFGFTdoc{Dok.~380 --- Warum $D_4$ --- Spezifität des Trägers gegen angepasste Vergleichsgitter (R95)}{../ch/380_D4_Spezifitaet_De_ch}
	\FFGFTdoc{Dok.~381 --- Die laufende Rekursion exakt aufsummiert}{../ch/381_Rekursion_Summation_De_ch}
	\FFGFTdoc{Dok.~382 --- Myon $g-2$: Stand 2025 und Einordnung der FFGFT-Aussagen}{../ch/382_Myon_g2_Stand2025_De_ch}
	\FFGFTdoc{Dok.~383 --- Massen und Planck-Skala ohne $v$}{../ch/383_Ein_Anker_ohne_v_De_ch}
	\FFGFTdoc{Dok.~384 --- FFGFT in Kurzfassung --- Grundlagen, Ergebnisse und Status nach der Rechenprüfung}{../ch/384_FFGFT_Kurzfassung_De_ch}
	\FFGFTdoc{Dok.~385 --- Der Higgs-Vakuum-Weg zu $\xi$}{../ch/385_Higgs_Vakuum_Weg_De_ch}
	\FFGFTdoc{Dok.~386 --- Wo $\alpha$ steckt --- Ladungseinheit, zwei Kugeln und $\xi$ als Fläche}{../ch/386_Alpha_Geometrie_De_ch}
	\FFGFTdoc{Dok.~387 --- Die Abweichungen im Überblick}{../ch/387_Abweichungen_De_ch}
	\FFGFTdoc{Dok.~388 --- Die CMB-Temperatur in der Grundform}{../ch/388_CMB_Grundform_De_ch}
	\FFGFTdoc{Dok.~389 --- Parabelflug, Trägheit und die Darstellungen der Gravitation}{../ch/389_Parabelflug_Darstellungen_De_ch}
\end{document}
